\section{Lecture 9 (September 22nd)}
\begin{prop}
(Two identical particles) Observe that the center of mass momentum and variables are
\[\vec{P}=\dfrac{m_2\vec{p}_{1}+m_1\vec{p}_{2}}{m_1+m_2}=\dfrac{1}{2}(\vec{p}_{1}-\vec{p}_{2})\qquad\&\qquad \vec{r}=\vec{r}_{1}-\vec{r}_{2}\qquad\&\qquad \mu=\dfrac{m_1m_2}{m_1+m_2}\]
For identical particles, the reduced mass becomes $\mu=m/2$, and in the center of mass coordinates, the Hamiltonian becomes
\[\left( E - \frac{\vec{P}_{\text{tot}}^{2}}{2M} \right)\psi(\vec{r})
= \left( \frac{\vec{p}^{2}}{2\mu} + V(|\vec{r}|) \right)\psi(\vec{r})\]
For two identical fermions, we observe the following constraints for $\psi ({\bf r})$:
\begin{itemize}
\item[(i)] Under interchange,
\[\psi (\vec{r})\mapsto -\psi (\vec{r})\]
\item[(ii)] And excluding spin and only the spatial part,
\[\psi \sim R_{nl}(r)Y_{lm}(\theta ,\phi )\]
with $Y_{lm}$ having $(-1)^{l}$ parity.
\item[(iii)] 
\[\dfrac{1}{2}\otimes \dfrac{1}{2}=\mathrm{\ triplet\ symmetry\ +\ singlet\ symmetry}\]
\end{itemize}
Therefore, since the total wavefunction must be antisymmetric, spin singlet (antimetric) $\times$ spatial symmetric is allowed requiring even $l$ ($0,2,4,\ldots $). Also, spin triplet (symmetric) $\times $ spatial antisymmetric is allowed requiring odd $l$ ($1,3,5,\ldots $).
\end{prop}
\vspace{2ex}
\begin{rmk}
(Intrinsic parity of fundamental particles) 
\begin{itemize}
\item[(i)] Let us first define the intrinsic parities of the electron, proton, and neutron as $+1$. 
\item[(ii)] Let us find the parity of the hydrogen atom at $l=1$. We find $(+1)(+1)(-1)=(-1)$. 
\item[(iii)] The parity of particles and anti-particles are opposite. Notice that for the positronium, $(+1)(-1)(+1)=0$.
\item[(iv)] As we know well the properties of the proton, neutron, and electron, we wanted to guess the properties of particles such as the pions ($\pi ^{-}$, $\pi ^{+}$, and $\pi ^{0}$). Also, we wanted to know particles like the deutron ($d^{+}=p^{+}-n$) and deuterium ($D=d^{+}-e^{-}$). After multiple particle experiments, we've found the capture reaction channel
\[d^{+}+\pi ^{-}\rightarrow n+n\]
and we have found that the total angular momentum of the initial system was 1 with the deuterium particle having spin 1. With the result being two identical fermions, we use the results above to analyse the final system. We have two possibilities, a spin singlet ($S=0$) or spin triplet ($S=1$). From spin conservation, we are able to know the specific state that the final state should be in.
\end{itemize}
The critical point of the last remark is that parity is conserved. In the last state, the parity is $-1$, and the inital state, we had 1 times the parity of $\pi ^{-1}$ telling us that the parity of $\pi ^{-1}$ is $-1$.
\end{rmk}
\vspace{2ex}
\begin{thm}
(Pauli exclusion principle of non-interacting fermions) We look at how the Pauli exclusion principle works for electrons. Consider a electron confined in a box of length $L$. The wavefunction of a single particle and its energy can be written as (through multiplying all components for the solution of an infinite potential well)
\[\Big(\dfrac{2}{L}\Big)^{3/2}\sin \Big(\dfrac{n_1\pi x}{L}\Big)\sin \Big(\dfrac{n_2\pi y}{L}\Big)\sin \Big(\dfrac{n_3\pi z}{L}\Big)\]
where $n_1,n_2,n_3\geq 1$. The energy should be given as
\[\dfrac{\hbar ^2}{2m}\Big[\Big(\dfrac{\pi n_1}{L}\Big)^2+\Big(\dfrac{\pi n_2}{L}\Big)^2+\Big(\dfrac{\pi n_3}{L}\Big)^2\Big]\]
and is independent of spin. Assume that there are 24 electrons inside the box. We want to see that the minimum energy is for these particles. In the lowest energy level of $3$, there are 2 possible particles and for the second lowest energy level of 6, there are 6. Also,
\[E(N=24)=214\times \Big(\dfrac{\hbar ^2\pi ^2}{2mL^2}\Big)\]
for fermions with half-spin. What about for bosons? We would have all particles stuffed in the lowest state and would have
\[E(N=24)=72\times \Big(\dfrac{\hbar ^2\pi ^2}{2mL^2}\Big)\]
What is the number of states with an energy level lower than $E$? It is
\[\dfrac{1}{8}\dfrac{\mathrm{Vol}}{\mathrm{\Delta }n_{x}\mathrm{\Delta }n_{y}\mathrm{\Delta }n_{z}}\]

\end{thm}
\vspace{2ex}

